% !TeX root = ../master.tex

\chapter{Conclusion}
\label{ch:conclusion}

\section{Summary}
\label{sec:summary}

% - Recap: investigated whether decoder-based LLMs overcome shortcut learning in AM
% - Three-part study: zero-shot robustness, context utilization, fine-tuning effects
% - Key finding 1: LLMs show substantially larger delta than encoders (RQ1)
% - Key finding 2: dataset-specific context improves zero-shot identification (RQ2)
% - Key finding 3: fine-tuning effects on shortcut learning (RQ3, TBD)
% - Conclusion on architecture-dependent vs. training-dependent shortcut learning

\section{Contributions}
\label{sec:final_contributions}

% - First systematic encoder vs. LLM shortcut learning comparison in argument mining
% - Word-shuffle manipulation as new diagnostic tool
% - Context utilization quantification across models and datasets
% - GAIC 2026 shared task submission
% - Reproducible experimental codebase

\section{Future Work}
\label{sec:future_work}

% - Chain-of-thought prompting as diagnostic for shortcut mitigation
% - Few-shot learning: does in-context learning introduce shortcuts like fine-tuning?
% - Multilingual evaluation: do findings transfer across languages?
% - Full AM pipeline: extend beyond identification to component + relation
% - Larger sample sizes and statistical significance testing
% - Other manipulation types: synonym replacement, paraphrase, back-translation
% - Fine-grained analysis: which specific linguistic features do LLMs attend to?
